\section{Regular rings}
\begin{theorem}[Auslander--Buchsbaum formula]\label{mcrt-thm-19-1}\dcref{19.1}\source{matsumura-commutative-ring-theory:page-0170}\uses{mcrt-thm-16-6}
If $A$ is Noetherian local and $0\ne M$ is finite with finite projective
dimension, then
\[
\operatorname{projdim}_A M+\depth M=\depth A.
\]
\end{theorem}
\begin{proof}
Induct on the projective dimension.  For dimension one, a minimal free
resolution yields a long exact Ext sequence because its matrix entries lie in
the maximal ideal; the general case follows by taking the first syzygy.
\end{proof}

\begin{theorem}[Serre's regularity criterion]\label{mcrt-thm-19-2}\dcref{19.2}\source{matsumura-commutative-ring-theory:page-0171}\uses{mcrt-thm-19-1,mcrt-def-koszul}
For a Noetherian local ring $A$,
\[
A\text{ regular}\quad\Longleftrightarrow\quad
\operatorname{gl.dim}A=\dim A\quad\Longleftrightarrow\quad
\operatorname{gl.dim}A<\infty.
\]
\end{theorem}
\begin{proof}
For a regular system of parameters the Koszul complex is a minimal resolution
of the residue field, giving global dimension equal to the dimension.  For the
converse, induct on embedding dimension, choose a regular element outside the
associated primes, split the induced syzygy map, and apply induction to the
quotient.
\end{proof}

\begin{theorem}[Localization of a regular ring]\label{mcrt-thm-19-3}\dcref{19.3}\source{matsumura-commutative-ring-theory:page-0172}\uses{mcrt-thm-19-2}
If $A$ is a regular local ring and $P$ is prime, then $A_P$ is regular.
\end{theorem}
\begin{proof}
A finite projective resolution of $A/P$ localizes to one of $K(P)$ over $A_P$;
Serre's criterion then applies.
\end{proof}

\begin{theorem}[Regular rings are normal]\label{mcrt-thm-19-4}\dcref{19.4}\source{matsumura-commutative-ring-theory:page-0172}\uses{mcrt-thm-19-3,mcrt-thm-17-8}
Every regular ring is normal.
\end{theorem}
\begin{proof}
Normality is local.  Height-one localizations are DVRs, and principal ideals
have no embedded prime divisors because regular local rings are CM; the normality
criterion for Noetherian domains follows.
\end{proof}

\begin{theorem}[Polynomial regularity]\label{mcrt-thm-19-5}\dcref{19.5}\source{matsumura-commutative-ring-theory:page-0172}\uses{mcrt-thm-19-2,mcrt-thm-19-3}
If $A$ is regular then $A[X]$ and $A[[X]]$ are regular.
\end{theorem}
\begin{proof}
Localize at maximal ideals, lift a monic residue polynomial, and compare height
and embedding dimension.  For power series use completion and invariance of
regularity under completion.
\end{proof}

\begin{theorem}[Free finite resolutions over depth-zero rings]\label{mcrt-thm-19-6}\dcref{19.6}\source{matsumura-commutative-ring-theory:page-0174}\uses{}
If every finite subset of the maximal ideal of a local ring is annihilated by a
nonzero element, then every module with a finite free resolution is free.
\end{theorem}
\begin{proof}
Take a minimal resolution and use Schanuel to see that the last syzygy is finite.
The annihilator hypothesis kills this free syzygy, so induction shortens the
resolution to length zero.
\end{proof}

\begin{theorem}[Euler characteristic positivity]\label{mcrt-thm-19-7}\dcref{19.7}\source{matsumura-commutative-ring-theory:page-0174}\uses{mcrt-thm-19-6}
If $M$ has a finite free resolution over a ring, then its Euler number is
nonnegative, and vanishing forces $M=0$ after localization at every minimal
prime.
\end{theorem}
\begin{proof}
Localize at a minimal prime and apply the preceding theorem to the resulting
zero-dimensional local ring; the Euler number is the rank of the free module.
\end{proof}

\begin{theorem}[Euler characteristic and regular elements]\label{mcrt-thm-19-8}\dcref{19.8}\source{matsumura-commutative-ring-theory:page-0175}\uses{mcrt-thm-19-7}
For a Noetherian ring and a module $M$ with a finite free resolution, the
following are equivalent: $\operatorname{Ann}(M)\ne0$; the Euler number
$\chi(M)=0$; and $\operatorname{Ann}(M)$ contains a regular element.
\end{theorem}
\begin{proof}
Localize at the associated primes and use Euler characteristic positivity.  The
annihilator then avoids every associated prime exactly when it contains a
regular element, yielding all three implications.
\end{proof}

\begin{theorem}[Vasconcelos' theorem]\label{mcrt-thm-19-9}\dcref{19.9}\source{matsumura-commutative-ring-theory:page-0175}\uses{mcrt-thm-19-8,mcrt-thm-16-2}
For a proper ideal $I$ in a Noetherian local ring, if
$\operatorname{projdim}_A I<\infty$, then $I$ is generated by an $A$-sequence
if and only if $I/I^2$ is free over $A/I$.
\end{theorem}
\begin{proof}
One implication follows from the graded pieces of a regular sequence.  For the
converse, choose a basis of $I/I^2$, split off a regular element using the
finite-resolution hypothesis, and induct on the number of generators.
\end{proof}
